\phantomsection
\addcontentsline{toc}{section}{References}
\begin{thebibliography}{99}
\raggedright

\bibitem{BBHS1998}
V. Bach, J.-M. Barbaroux, B. Helffer, and H. Siedentop,
\emph{Stability of Matter for the Hartree--Fock Functional of the Relativistic Electron-Positron Field},
Documenta Mathematica \textbf{3} (1998), 353--364.
\href{https://doi.org/10.4171/DM/51}{doi:10.4171/DM/51}.
\href{https://ems.press/journals/dm/articles/8964982}{Journal page}.

\bibitem{ELS2021}
M.~J. Esteban, M. Lewin, and \'{E}. S\'er\'e,
\emph{Dirac--Coulomb Operators with General Charge Distribution. I. Distinguished Extension and Min-Max Formulas},
Annales Henri Lebesgue \textbf{4} (2021), 1421--1456.
\href{https://doi.org/10.5802/ahl.106}{doi:10.5802/ahl.106}.
\href{https://arxiv.org/abs/2003.04004}{arXiv:2003.04004}.

\bibitem{FSS}
J. Fr\"ohlich, B. Simon, and T. Spencer,
\emph{Infrared bounds, phase transitions and continuous symmetry breaking},
Communications in Mathematical Physics \textbf{50} (1976), 79--95.
\href{https://doi.org/10.1007/BF01608557}{doi:10.1007/BF01608557}.

\bibitem{GHLS2013}
P. Gravejat, C. Hainzl, M. Lewin, and \'E. S\'er\'e,
\emph{Construction of the Pauli--Villars-Regulated Dirac Vacuum in Electromagnetic Fields},
Archive for Rational Mechanics and Analysis \textbf{208} (2013), 603--665.
\href{https://doi.org/10.1007/s00205-012-0609-1}{doi:10.1007/s00205-012-0609-1}.
\href{https://arxiv.org/abs/1204.2893}{arXiv:1204.2893}.

\bibitem{GLS2009}
P. Gravejat, M. Lewin, and \'{E}. S\'er\'e,
\emph{Ground State and Charge Renormalization in a Nonlinear Model of Relativistic Atoms},
Communications in Mathematical Physics \textbf{286} (2009), 179--215.
\href{https://doi.org/10.1007/s00220-008-0660-9}{doi:10.1007/s00220-008-0660-9}.
\href{https://arxiv.org/abs/0712.2911}{arXiv:0712.2911}.

\bibitem{HLS2004}
C. Hainzl, M. Lewin, and \'{E}. S\'er\'e,
\emph{Existence of a Stable Polarized Vacuum in the Bogoliubov--Dirac--Fock Approximation},
Communications in Mathematical Physics \textbf{257} (2005), 515--562.
\href{https://doi.org/10.1007/s00220-005-1343-4}{doi:10.1007/s00220-005-1343-4}.
\href{https://arxiv.org/abs/math-ph/0403005}{arXiv:math-ph/0403005}
(2004 preprint).

\bibitem{HLS2005}
C. Hainzl, M. Lewin, and \'{E}. S\'er\'e,
\emph{Self-consistent Solution for the Polarized Vacuum in a No-photon QED Model},
Journal of Physics A: Mathematical and General \textbf{38} (2005), 4483--4499.
\href{https://doi.org/10.1088/0305-4470/38/20/014}{doi:10.1088/0305-4470/38/20/014}.
\href{https://arxiv.org/abs/physics/0404047}{arXiv:physics/0404047}.

\bibitem{HLS2008}
C. Hainzl, M. Lewin, and \'{E}. S\'er\'e,
\emph{Existence of Atoms and Molecules in the Mean-Field Approximation of No-Photon Quantum Electrodynamics},
Archive for Rational Mechanics and Analysis \textbf{192} (2009), 453--499.
\href{https://doi.org/10.1007/s00205-008-0144-2}{doi:10.1007/s00205-008-0144-2}.
\href{https://arxiv.org/abs/math-ph/0606001}{arXiv:math-ph/0606001}
(final preprint version 2008).

\bibitem{HLSol2007}
C. Hainzl, M. Lewin, and J.~P. Solovej,
\emph{The Mean-Field Approximation in Quantum Electrodynamics: the No-Photon Case},
Communications on Pure and Applied Mathematics \textbf{60} (2007), 546--596.
\href{https://doi.org/10.1002/cpa.20145}{doi:10.1002/cpa.20145}.
\href{https://arxiv.org/abs/math-ph/0503075}{arXiv:math-ph/0503075}.

\bibitem{JW2018}
C.~D. J\"akel and W.~F. Wreszinski,
\emph{Stability of Relativistic Quantum Electrodynamics in the Coulomb Gauge},
Journal of Mathematical Physics \textbf{59} (2018), 032303.
\href{https://doi.org/10.1063/1.5011031}{doi:10.1063/1.5011031}.
\href{https://arxiv.org/abs/1705.04652}{arXiv:1705.04652}.
The 2017 preprint and 2018 published version are distinguished where relevant.

\bibitem{Lieb2002}
E.~H. Lieb,
\emph{The Stability of Matter and Quantum Electrodynamics},
Heisenberg Symposium contribution (2002).
\href{https://arxiv.org/abs/math-ph/0209034}{arXiv:math-ph/0209034}.

\bibitem{LiJiangYao}
Z.-X. Li, Y.-F. Jiang, and H. Yao,
\emph{Majorana-Time-Reversal Symmetries: A Fundamental Principle for
Sign-Problem-Free Quantum Monte Carlo Simulations},
Physical Review Letters \textbf{117} (2016), 267002.
\href{https://doi.org/10.1103/PhysRevLett.117.267002}{doi:10.1103/PhysRevLett.117.267002};
\href{https://arxiv.org/abs/1601.05780}{arXiv:1601.05780}.

\bibitem{LL2002}
E.~H. Lieb and M. Loss,
\emph{Stability of a Model of Relativistic Quantum Electrodynamics},
Communications in Mathematical Physics \textbf{228} (2002), 561--588.
\href{https://doi.org/10.1007/s002200200665}{doi:10.1007/s002200200665}.
\href{https://arxiv.org/abs/math-ph/0109002}{arXiv:math-ph/0109002}.

\bibitem{LS2000}
E.~H. Lieb and H. Siedentop,
\emph{Renormalization of the Regularized Relativistic Electron-Positron Field},
Communications in Mathematical Physics \textbf{213} (2000), 673--683.
\href{https://doi.org/10.1007/s002200000265}{Publisher DOI}.
\href{https://arxiv.org/abs/math-ph/0003001}{arXiv:math-ph/0003001}.

\bibitem{LY1988}
E.~H. Lieb and H.-T. Yau,
\emph{The Stability and Instability of Relativistic Matter},
Communications in Mathematical Physics \textbf{118} (1988), 177--213.
\href{https://doi.org/10.1007/BF01218577}{doi:10.1007/BF01218577}.
\href{https://collaborate.princeton.edu/en/publications/the-stability-and-instability-of-relativistic-matter-2/}{Author's institutional publication record}.

\bibitem{Miyao}
T. Miyao, \emph{Upper bounds on the charge susceptibility of many-electron
systems coupled to the quantized radiation field},
Letters in Mathematical Physics \textbf{105} (2015), 1119--1133.
\href{https://doi.org/10.1007/s11005-015-0775-9}{doi:10.1007/s11005-015-0775-9};
\href{https://arxiv.org/abs/1405.0112}{arXiv:1405.0112}.

\bibitem{VafaWitten}
C. Vafa and E. Witten,
\emph{Restrictions on symmetry breaking in vector-like gauge theories},
Nuclear Physics B \textbf{234} (1984), 173--188.
\href{https://doi.org/10.1016/0550-3213(84)90230-X}{doi:10.1016/0550-3213(84)90230-X}.

\bibitem{Wei}
Z. C. Wei, C. Wu, Y. Li, S. Zhang, and T. Xiang,
\emph{Majorana Positivity and the Fermion Sign Problem of Quantum Monte Carlo Simulations},
Physical Review Letters \textbf{116} (2016), 250601.
\href{https://doi.org/10.1103/PhysRevLett.116.250601}{doi:10.1103/PhysRevLett.116.250601};
\href{https://arxiv.org/abs/1601.01994}{arXiv:1601.01994}.
\bibitem{Herbst1977}
I. W. Herbst,
\emph{Spectral Theory of the Operator \((p^2+m^2)^{1/2}-Ze^2/r\)},
Communications in Mathematical Physics \textbf{53} (1977), 285--294.
\href{https://doi.org/10.1007/BF01609852}{doi:10.1007/BF01609852}.
\end{thebibliography}
